\section{Condensed three-level closure}
\label{sec:scalar15handoff}

The P3 spine fixes the common assembly language but does not itself choose a
non-Abelian particle completion.  A later P4 branch closes three previously
separate questions at three different logical levels.  We record that result
here because it is the first particle-sector specialization of the joint
source architecture; the underlying P4 reconstruction is not reclassified as
a theorem owned by Paper VII.

\subsection{Class-relative base--seed ownership}

Let \(\mathcal C_{15}^{\rm rec}\) denote the recovery-complete,
source-generated and source-irredundant P4/type-II class.  It contains the
recovered one-copy PIS--CAR packet, one occupied Higgs branch, the scalar-15
terminal and no independently sourced terminal-null operation ideal.  The
reverse source reconstruction identifies the pointed exterior \(V_5\) seed
grammar and a positive Hermitian/Hodge base realization.  The one-carrier
type-II selector then fixes the typed edge
\begin{equation}
 e_{HH\Delta}:\operatorname{Sym}^2V_5\longrightarrow S_{15},
 \qquad \dim\operatorname{Sym}^2V_5=15,
 \label{eq:scalar15edge}
\end{equation}
and faithful realization gives rank fifteen.  The unique shortest connected
alternating seed--Higgs word and its positive quartic completion imply
\begin{equation}
 \boxed{
 \mathcal C_{15}^{\rm rec}
 \Longleftrightarrow
 (\mathcal S_{\wedge,5},e_{HH\Delta};h_B,*_B)
 \Longrightarrow h_*\ne0,\quad\Delta_*\ne0 .}
 \label{eq:scalar15ownership}
\end{equation}
Thus charged character, occupation and quartic stabilization require no
independent source datum \emph{inside this recovery class}.  Equation
\eqref{eq:scalar15ownership} is a minimal-preimage reconstruction of the
selected branch.  It is not a theorem that the unrestricted physical parent
occupies type II rather than another neutrino/particle completion.

\subsection{Common spectral scale}

The same completion relates its Higgs length and scalar-15 mass to the common
Planck terminal:
\begin{equation}
 \ell_H=\sqrt{48\pi}\,\ell_P,
 \qquad
 M_*=\frac{\sqrt2\,\hbar c}{\ell_H}
 =\frac{E_P}{\sqrt{24\pi}}
 =\frac{\overline M_Pc^2}{\sqrt3},
 \label{eq:scalar15scale}
\end{equation}
where \(E_P=\sqrt{\hbar c^5/G}\) and
\(\overline M_Pc^2=E_P/\sqrt{8\pi}\).  With the 2022 CODATA constants
\citep{CODATA2022},
\begin{equation}
 \boxed{M_*=1.406035(16)\times10^{18}\ {\rm GeV}.}
 \label{eq:scalar15numeric}
\end{equation}
The Tau-specific statement is the dimensionless ratio
\(M_*/E_P=1/\sqrt{24\pi}\) and the absence of an independent scalar-15 length.
Equation~\eqref{eq:scalar15numeric} is a one-anchor metrological calibration,
not a first-principles prediction of the SI value of \(G\).

\subsection{Exact interacting functional law}

Lift the finite particle operator to the full graded four-dimensional carrier
and let \(P\) be the positive background-covariant Laplace-type operator,
including the declared gauge-fixing and ghost blocks.  The no-extra-profile
spectral scheme uses
\begin{equation}
 \boxed{R_k(P)=\frac{P}{\exp(P/k^2)-1}}
 \label{eq:scalar15regulator}
\end{equation}
in the exact effective-average-action flow \citep{Wetterich1993}
\begin{equation}
 \boxed{
 \partial_t\Gamma_k=\frac{1}{2}\operatorname{STr}
 \left[(\Gamma_k^{(2)}+R_k)^{-1}\partial_tR_k\right],
 \qquad t=\log k .}
 \label{eq:scalar15flow}
\end{equation}
For gauge sectors Eq.~\eqref{eq:scalar15flow} is solved together with the
regulator-modified Slavnov--Taylor/BV identity; the regulator insertion must
vanish as \(k\to0\).  If \(\{\mathcal O_I\}\) is a radiatively and
BRST-complete independent operator basis and
\(G_{IJ}=\langle\mathcal O_I,\mathcal O_J\rangle\) on the source-owned trace
domain, the component beta functions are the unique projection
\begin{equation}
 \boxed{
 \beta^I=(G^{-1})^{IJ}
 \left\langle\mathcal O_J,\partial_t\Gamma_k\right\rangle .}
 \label{eq:scalar15betaprojection}
\end{equation}
Equations~\eqref{eq:scalar15regulator}--
\eqref{eq:scalar15betaprojection} fix one gauge-covariant exact functional
law, threshold prescription and coordinate projection.  Computing the full
component polynomials remains an evaluation problem.  The regulator is an
analysis scheme, not a physical coupling, observer quantizer or additional
law requiring Nature-level selection; exact \(k=0\) observables are regulator
independent, whereas residual dependence of a finite truncation is an
uncertainty diagnostic.

The net closure is therefore deliberately three-tiered:
\begin{center}
\begin{tabular}{@{}ll@{}}
\toprule
Question & Status\\
\midrule
base--seed ownership & proved inside \(\mathcal C_{15}^{\rm rec}\)\\
absolute numerical scale & dimensionless ratio derived; SI value calibrated once\\
full beta law & exact functional flow and projection fixed\\
physical type-II occupation & empirical model-selection question\\
\bottomrule
\end{tabular}
\end{center}
